\section*{Chapter \ref{ch:b3:gfd-foundations} --- Geophysical Fluid Dynamics}

\begin{solution}{exo:b3:gfd-foundations:1}
(a)~$f = 2\Omega\sin\varphi$: \num{7.292e-5}, \num{1.031e-4} and
\qty{1.263e-4}{s^{-1}}; $2\pi/f$: 23.93, 16.92 and \qty{13.82}{h}.
(b)~$U/f$: 2.74, 1.94 and \qty{1.58}{km}. (c)~At \ang{30}, $\sin\varphi =
1/2$ and $2\pi/f = 2\pi/\Omega$: the inertial period is the sidereal day,
\qty{23.93}{h}.
\end{solution}

\begin{solution}{exo:b3:gfd-foundations:2}
With $f = \qty{1.031e-4}{s^{-1}}$: dust devil \num{9.7e3}, sea breeze 2.4,
depression 0.15, gyre \num{2.4e-4}. The depression and the gyre are close to
geostrophic balance; the dust devil turns either way, unaware of the
Earth's rotation, and the sea breeze feels it only over a day.
\end{solution}

\begin{solution}{exo:b3:gfd-foundations:3}
$\partial u_g/\partial z = (g/fT)\,|\partial T/\partial y| = 9.81\times
\num{5e-6}/(\num{1.031e-4}\times250) = \qty{1.90e-3}{s^{-1}}$; over
\qty{10}{km} the westerly wind increases by \qty{19}{m/s}.
\end{solution}

\begin{solution}{exo:b3:gfd-foundations:4}
$E = \num{50e-8}$: $N = \qty{2.21e-3}{s^{-1}}$, a period of \qty{47}{min},
\num{1.27} cycles per hour. $E = \num{1000e-8}$: $N = \qty{9.90e-3}{s^{-1}}$,
\qty{10.6}{min}, 5.7 cycles per hour. $E = \num{e-8}$: $N =
\qty{3.13e-4}{s^{-1}}$, \qty{5.6}{h}, 0.18 cycles per hour.
\end{solution}

\begin{solution}{exo:b3:gfd-foundations:5}
(a)~$W \sim UH/L = \qty{0.1}{m/s}$ and $\mathrm Dw/\mathrm Dt \sim UW/L =
\qty{e-6}{m/s^2}$, \num{1.0e-7} of $g$. The pressure term, $fUL/H =
\qty{0.103}{m/s^2}$, is larger by \num{1e5}: the ratio, \num{9.7e-6}, is
$\mathrm{Ro}(H/L)^2$ with $\mathrm{Ro} = 0.097$. (b)~$\tilde fu =
2\Omega\cos\ang{45}\times10 = \qty{1.03e-3}{m/s^2}$, \num{1.05e-4} of $g$.
\end{solution}

\begin{solution}{exo:b3:gfd-foundations:6}
$fr = \qty{30}{m/s}$. (a)~$\mathrm{Ro}_g = 0.667$, $V = 40/(1 + \sqrt{3.667})
= \qty{13.7}{m/s}$, a third slower than geostrophy. (b)~$\mathrm{Ro}_g =
0.167$, $V = 10/(1 + \sqrt{0.333}) = \qty{6.34}{m/s}$. (c)~$V_g \le fr/4 =
\qty{7.5}{m/s}$, and then $V = fr/2 = \qty{15}{m/s}$, twice the geostrophic
speed.
\end{solution}

\begin{solution}{exo:b3:gfd-foundations:7}
(a)~$\Omega = 2\pi/6 = \qty{1.047}{rad/s}$, $f = \qty{2.094}{s^{-1}}$.
$\mathrm{Ro} = 0.01/(2.094\times0.05) = 0.095$ and $\mathrm{Ek} =
\num{e-6}/(2.094\times0.04) = \num{1.19e-5}$: both small, so the flow obeys
the Taylor--Proudman theorem and a \omterm{thm:b3:gfd-foundations:taylor-proudman}{Taylor column} forms. (b)~$0.2/\sqrt{\num{e-6}
\times1.047} = \qty{195}{s}$, \qty{3.3}{min}, 33 turns of the table.
\end{solution}

\begin{solution}{exo:b3:gfd-foundations:8}
$\partial\bar T/\partial y = -f\,\Delta u/(R\ln(850/250))$, with
$\ln(850/250) = 1.224$. At \ang{30} N, $-\num{7.29e-5}\times36.12/(286.9
\times1.224) = \qty{-7.50}{K}$ per \qty{1000}{km}; at \ang{45} N,
$-\num{1.031e-4}\times17.18/351.1 = \qty{-5.05}{K}$ per \qty{1000}{km}. The
shear is proportional to the gradient divided by $f$: the gradient is 1.5
times larger at \ang{30} and $f$ 1.4 times smaller, so the shear is 2.1 times
larger.
\end{solution}

\begin{solution}{exo:b3:gfd-foundations:9}
(a)~$L_R = 0.0113\times11\,000/\num{1.031e-4} = \qty{1206}{km}$;
$0.00357\times1000/\num{9.47e-5} = \qty{37.7}{km}$, 32 times smaller.
(b)~$\beta = 2\Omega\cos\varphi/a$: \num{1.62e-11} and
\qty{1.74e-11}{m^{-1}.s^{-1}}. $\beta L_R/f = 0.19$ for the atmosphere, 0.007
for the ocean.
\end{solution}

\begin{solution}{exo:b3:gfd-foundations:10}
(a)~$\mathrm dw/\mathrm dt = -(r + \mathrm if)w$, so $w = U\mathrm
e^{-rt}\mathrm e^{-\mathrm ift}$: the velocity turns at the rate $f$ while
its magnitude decays as $\mathrm e^{-rt}$. The position $x + \mathrm iy =
U(1 - \mathrm e^{-(r + \mathrm if)t})/(r + \mathrm if)$ spirals into the
point $U/(r + \mathrm if)$. (b)~$x + \mathrm iy = U(1 - \mathrm e^{-(r +
\mathrm if)t})/(r + \mathrm if) + u_0t$: the spiral is drawn out into loops.
(c)~$0.2\,\mathrm e^{-1} = \qty{0.074}{m/s}$; the centre has drifted $u_0t =
0.05\times4\times86\,400 = \qty{17.3}{km}$ eastward.
\end{solution}

\begin{solution}{exo:b3:gfd-foundations:11}
$f = \qty{9.37e-5}{s^{-1}}$; $\partial u/\partial z = (g/f\rho_0)\,\partial
\rho/\partial y = 9.81\times10^{-5}/(\num{9.37e-5}\times1025) =
\qty{1.02e-3}{s^{-1}}$. With $u = 0$ at \qty{1000}{m} the surface current is
\qty{1.02}{m/s}, eastward: in the Northern Hemisphere the current flows with
the light water on its right.
\end{solution}

\begin{solution}{exo:b3:gfd-foundations:12}
(a)~$N^2 = (g/\rho)\,\Delta\sigma_\theta/\Delta z$. 0--\qty{100}{m}:
$9.81\times0.668/(1026.4\times100)$, $N = \qty{7.99e-3}{s^{-1}}$, 4.6
cycles per hour; 100--\qty{500}{m}: \qty{3.02e-3}{s^{-1}}, 1.7 cycles per
hour; 500--\qty{1000}{m}: \qty{3.24e-3}{s^{-1}}, 1.9 cycles per hour.
(b)~The top \qty{100}{m}: it contains the warm surface water above the
seasonal pycnocline. Between 100 and \qty{500}{m} the water is more
uniform than below. (c)~$E = N^2/g$ is \num{650e-8}, \num{93e-8} and
\qty{107e-8}{m^{-1}}, all within the range of the upper ocean.
\end{solution}

\begin{solution}{pb:b3:gfd-foundations:1}
\textbf{1.} $f = 2\times\num{7.292115e-5}\sin\ang{48} =
\qty{1.084e-4}{s^{-1}}$; $2\pi/f = \qty{16.1}{h}$.

\textbf{2.} Bathtub \num{3.1e3}; tank ($f = 4\pi/10 = \qty{1.257}{s^{-1}}$)
0.020; breeze 2.3; depression 0.14; gyre \num{2.3e-4}.

\textbf{3.} The bathtub feels nothing of the Earth's rotation; the breeze
feels it only after some hours. The tank, the depression and the gyre have
small \omterm{def:b3:gfd-foundations:numbers}{Rossby numbers} and are close to geostrophic.

\textbf{4.} Tank $\num{e-6}/(1.257\times0.04) = \num{2.0e-5}$; depression
$10/(\num{1.084e-4}\times10^8) = \num{9.2e-4}$; gyre
$\num{e-3}/(\num{1.084e-4}\times10^6) = \num{9.2e-6}$.

\textbf{5.} $\delta/H = \sqrt{2\nu/(fH^2)} = \sqrt{2\,\mathrm{Ek}}$. Tank
$\sqrt{2\times\num{e-6}/1.257} = \qty{1.3}{mm}$; air $\sqrt{20/\num{1.084e-4}}
= \qty{430}{m}$.

\textbf{6.} Friction acts in layers a few thousandths of the depth in the
tank and the ocean, a few hundredths in the atmosphere; above them the flow
is nearly frictionless.

\textbf{7.} $\rho = p/(RT) = 101\,290/(286.9\times288.1) =
\qty{1.225}{kg/m^3}$.

\textbf{8.} $u_g = (1/\rho f)\,|\partial p/\partial y| = (400/300\,000)/(1.225
\times\num{1.084e-4}) = \qty{10.0}{m/s}$, from the west, with the low
pressure to the north on the left.

\textbf{9.} $\mathrm{Ro}_g = 10.04/(\num{1.084e-4}\times600\,000) = 0.154$;
$V = 2\times10.04/(1 + \sqrt{1.618}) = \qty{8.84}{m/s}$, \qty{12}{\percent}
slower.

\textbf{10.} $V = 2\times10.04/(1 + \sqrt{0.382}) = \qty{12.4}{m/s}$. The
largest $V_g$ is $fr/4 = \qty{16.3}{m/s}$, a gradient $\rho fV_g =
\qty{2.16e-3}{Pa/m}$, \qty{6.5}{hPa} per \qty{300}{km}.

\textbf{11.} $g/c_p = \qty{9.77}{K/km}$. $N^2 = (9.81/255.7)\times(9.77 -
6.49)\times10^{-3} = \qty{1.26e-4}{s^{-2}}$, $N = \qty{0.0112}{s^{-1}}$, a
period of \qty{9.3}{min}.

\textbf{12.} With $\Gamma = 0$, $N^2 = (g/T)(g/c_p) = g^2/(c_pT) =
\qty{4.43e-4}{s^{-2}}$, $N = \qty{0.0210}{s^{-1}}$, 1.9 times the
tropospheric value.

\textbf{13.} $L_R = 0.0112\times10^4/\num{1.084e-4} = \qty{1035}{km}$, the
size of the depression.

\textbf{14.} $g' = \qty{0.0147}{m/s^2}$; $\sqrt{g'h} = \qty{0.54}{m/s}$;
$L_R = 0.542/\num{1.084e-4} = \qty{5.0}{km}$.

\textbf{15.} Yes: the lake, \qty{12}{km} wide, is more than twice its
Rossby radius, so its internal waves and currents are turned by the
Coriolis force.

\textbf{16.} $(249.4 - 253.0)/600 = \qty{-6.0}{K}$ per \qty{1000}{km}.

\textbf{17.} $\ln(1013/265) = 1.3409$. South:
$286.9\times253.0\times1.3409/9.81 = \qty{9922}{m}$; north \qty{9781}{m}.
The \qty{265}{hPa} surface lies \qty{141}{m} lower in the north, a slope of
\num{2.35e-4} relative to the lower surface.

\textbf{18.} $\Delta u = (g/f)\times\text{slope} = 9.81\times\num{2.35e-4}/
\num{1.084e-4} = \qty{21.3}{m/s}$.

\textbf{19.} $(R/f)\ln(p_b/p_t)\,|\partial\bar T/\partial y| =
(286.9/\num{1.084e-4})\times1.3409\times\num{6.0e-6} = \qty{21.3}{m/s}$, the
same: the two calculations are one.

\textbf{20.} \omterm{def:b3:gfd-foundations:baroclinic}{Baroclinic}: the temperature, hence the density, changes along
the pressure surfaces. \omterm{def:b3:gfd-foundations:baroclinic}{Barotropic}, the wind would not change with height:
\qty{10.0}{m/s}.

\textbf{21.} $(g/f\bar T)\,|\partial T/\partial y|\,\Delta z = 9.81\times
\num{6.0e-6}\times10^4/(\num{1.084e-4}\times255.7) = \qty{21.2}{m/s}$,
within \qty{0.3}{\percent}.

\textbf{22.} $10.04 + (286.9/\num{1.084e-4})\times\num{6.0e-6}\times
\ln(1013/541) = 10.04 + 9.96 = \qty{20.0}{m/s}$.

\textbf{23.} $\partial u/\partial z = 21.3/10^4 = \qty{2.13e-3}{s^{-1}}$;
$\mathrm{Ri} = \num{1.26e-4}/\num{4.54e-6} = 28$, far above 0.25: the shear
does not make the free troposphere turbulent.

\textbf{24.} The wind of Orrin, \qty{31}{m/s}, is stronger than the zonal mean
at \ang{45} N and weaker than at \ang{30} N. The polar-front jet wanders
north and south; at any moment its gradient, and its wind, is concentrated
over a band, and averaging round a latitude circle spreads it out.

\textbf{25.} $10.0 + 21.3$: the geostrophic wind at \qty{10}{km} above the
steppe is \textbf{\qty[text-series-to-math=true]{31}{m/s}}, from the west
(\qty{31.3}{m/s} by the calculation).
\end{solution}
